% \begin{savequote}[75mm]
% Quote
% \qauthor{Quoteauthor Lastname}
% \end{savequote}
% \newthought{There's something to be said} for having a good opening line. 

% Test citation 1 \citep{Azevedo2017-kj} and  2 \citet{Izhikevich2010-tq}.
% For an example of a full page figure, see Figure \ref{fig:fig_01}.

\chapter{General Introduction and Background}

Intelligent systems do not simply store information—they transform inputs into internal representations that can be acted upon. In neuroscience, a central question is how the brain represents sensory evidence and learning signals in a form that can be read out by motor circuits to generate adaptive behavior. In foundation models, an analogous question is how training objectives and internal representations shape what is learned and how effectively it supports downstream task performance. In both settings, learning is only useful to the extent that it produces representations that are accessible to the relevant readout—motor control in animals, and inference-time prompting/decoding interfaces in models—and robust across changes in context.

This dissertation examines that principle across three systems that share a common computational demand: converting sparse or structured cues into reliable outputs. The unifying focus is not the domain (odor, text, or images), but the coupling between (i) what is represented internally, (ii) how learning shapes those representations, and (iii) how those representations are ultimately read out to drive behavior or task success.

In the first part, I study how animals convert sparse sensory evidence into motor policy during olfactory navigation. Odor signals in natural environments are intermittent and fragmented, so successful navigation requires learning and representing not only instantaneous stimulus identity but also the statistics of recent encounters. Using a closed-loop olfactory navigation paradigm with visually anchored heading, I show that flies express distinct control modes at odor boundaries: when encounters are frequent, odor onset stabilizes heading and promotes forward locomotion; when encounters are sparse, the same sensory transition promotes turning and reorientation. I then link these behavioral outputs to specific neural representations—an opponent code in mushroom body output neurons that tracks encounter statistics and its coupling to heading representations in the central complex—providing a circuit-level account of how learned sensory history is transformed into action selection.

In the second part, I turn to a parallel question in generative vision models: what internal representations and mechanisms determine whether a model can satisfy structured constraints at inference time. In text-to-image diffusion, a particularly diagnostic constraint is two-object spatial relations, where “success” is whether generated images place objects in the correct relation (often despite strong fidelity on single-object attributes). Using controlled training and mechanistic analysis, I dissect how diffusion transformers implement relational generation and why this capability can be fragile. I identify a structured multi-step circuit in which relation information is converted into spatially localized intermediate signals and then used to assign object identity, providing a mechanistic explanation of when relational generation succeeds and when it breaks under perturbations. More broadly, this chapter frames downstream performance as a readout problem: relation information may exist in the model, but its utility depends on whether it is routed through a robust, readable pathway during generation.

In the third part, I study an explicit form of learning and representation update in language models: injecting new factual knowledge directly into model weights. This differs from conventional fine-tuning settings where training and evaluation share a tightly matched format (e.g., supervised examples and similarly structured tests). Here, the goal is to update a pretrained autoregressive model using only raw text describing new knowledge—without constructing question–answer supervision—and then evaluate whether the model can later express that knowledge in response to natural queries. In controlled knowledge-injection experiments, standard autoregressive fine-tuning often fails to generalize from documents to question answering without heavy paraphrase augmentation and can exhibit order-sensitive failures such as reversal-like effects. I test the hypothesis that these limitations reflect a mismatch between the autoregressive objective and the downstream readout interface. By reframing weight updates as a masked reconstruction objective inspired by demasking-style training, I show that changing the objective can substantially improve knowledge injection from raw text and reduce reversal-type failures, highlighting how objective design determines whether new information becomes usable at inference time.

Taken together, these projects support a shared conclusion across neuroscience and foundation models: learning shapes internal representations, but whether those representations improve behavior or downstream task performance depends on how effectively they can be read out under realistic conditions. By treating representation, objective, and readout as a coupled system, this dissertation connects circuit-level mechanisms of sensorimotor control with mechanistic constraints underlying compositional generation and knowledge updating in modern foundation models.

\afterpage{\clearpage}